\section{The Dimension Ladder}

The selection argument has two halves. This section is the first. It shows that only a four-dimensional bulk
yields the three-dimensional physical space we appear to inhabit. The argument turns on one geometric fact about
hyperbolic honeycombs, and that fact is sharp enough to rule out every dimension but one. The second half, which
singles out $\{3,4,3,4\}$ within the four-dimensional family, is taken up in the next section.

\subsection{Physical space is the bulk dimension minus one}

The flat layer that matter lives on is not the bulk itself. It is the cusp, the horosphere, the flat slice that
sits at the ideal boundary of the hyperbolic crystal. That slice is always exactly one dimension below the bulk.
A four-dimensional bulk has a three-dimensional cusp, a three-dimensional bulk has a two-dimensional cusp, and so
on down. This single relation fixes which substrate dimension can give the observed space.

\begin{proposition}[Dimension ladder]
The dimension of the physical space carried by a hyperbolic honeycomb equals the bulk dimension minus one. Among
the candidate families, only the four-dimensional honeycombs land on a three-dimensional physical space.
\end{proposition}

The table below walks the ladder. The two-dimensional and three-dimensional families undershoot, giving
one-dimensional and two-dimensional space. The five-dimensional family overshoots, giving four-dimensional space.
Only the four-dimensional family is exact.

\begin{table}[ht]
\centering
\begin{tabular}{llll}
\toprule
Substrate family & Bulk & Physical space (cusp) & Verdict \\
\midrule
$2$D tilings, e.g. $\{7,3\}$ & $2$D & $1$D & under \\
$3$D honeycombs, e.g. $\{5,3,4\}$ & $3$D & $2$D & under \\
$4$D honeycombs, $\{3,4,3,4\}$ & $4$D & $3$D & exact \\
$5$D honeycombs, e.g. $\{5,3,3,3,3\}$ & $5$D & $4$D & over \\
\bottomrule
\end{tabular}
\caption{The cusp is one dimension below the bulk, so only the four-dimensional family gives the observed three.}
\end{table}

\subsection{What the wrong dimension costs}

The dimension of the physical space is not an abstract bookkeeping number. It sets the kind of matter and the law
of gravity that the substrate can support, and each wrong dimension fails in a recognisable way. A
one-dimensional space, from a two-dimensional substrate, has no exchange statistics at all, a confining linear
gravity, and only kink solitons. A two-dimensional space, from a three-dimensional substrate, has anyonic
statistics, a logarithmic gravity, and two-dimensional skyrmions. A four-dimensional space, from a
five-dimensional substrate, has $1/r^2$ gravity, four-dimensional instantons, and one dimension too many. A
three-dimensional space, from a four-dimensional substrate, has fermions, $1/r$ gravity, and three-dimensional
hopfions. That last row is the observed case. The matter statistics and the gravity law both track the ladder,
and only the four-dimensional substrate gives the fermions and the $1/r$ gravity we measure.

\subsection{The dimension window}

One might worry that hyperbolic crystals are scarce, and that the ladder above is choosing among a handful of
oddities. The opposite is true. Regular hyperbolic honeycombs are plentiful, and counting them by compactness
class makes the real constraint visible. We use the taxonomy of the geometry section, where compact means every
element is bounded, paracompact means there is one ideal element at infinity, and noncompact means there is an
infinite ideal cell.

\begin{table}[ht]
\centering
\begin{tabularx}{\textwidth}{p{2.2cm}p{2.6cm}Xp{1.6cm}}
\toprule
Dimension & Compact regular & Paracompact or noncompact regular & Physical space \\
\midrule
$2$, $H^2$ & infinitely many, $\{7,3\}$, $\{5,4\}$, \dots & infinitely many more, apeirogonal, $\{\infty,3\}$, \dots & $1$D \\
$3$, $H^3$ & $4$, all contain a $5$ & $11$ more, the $3$/$4$/$6$ family, ideal cells or vertices & $2$D \\
$4$, $H^4$ & $5$, all contain a $5$ & several more, including $\{3,4,3,4\}$ with its ideal $\{4,3,4\}$ cubic cusp & $3$D \\
$5$, $H^5$ & $0$ & a few, the last dimension with any regular honeycomb, including the $D_4$ hook $\{3,4,3,3,4\}$ & $4$D \\
$6$ and up & $0$ & $0$ & --- \\
\bottomrule
\end{tabularx}
\caption{Regular hyperbolic honeycombs by compactness class. The window has two teeth.}
\end{table}

The window has two teeth, and together they cut the field down to one dimension. The first tooth is compactness.
The bounded regulars, the ones with no cusp, exist only in dimensions two through four, and every one of them
carries a $5$ in its symbol. That places them in the icosahedral family, which is non-crystallographic. They
cannot carry a root system, so they have no spinor and no gauge group. The second tooth is the
crystallographic restriction. The honeycombs whose symbols use only $3$s and $4$s, the ones that can carry a
root system, are all paracompact. The substrate $\{3,4,3,4\}$ is among them, and its ideal element is the flat
$\{4,3,4\}$ cubic cusp. Compactness in dimensions two through four forces a $5$ and kills the spinor. The
spinor-bearing crystallographic substrate must therefore be paracompact. Four dimensions is the only paracompact
crystallographic dimension that yields a three-dimensional cusp, since five overshoots to four and nothing
regular survives above five.

\subsection{An independent confirmation of three dimensions}

There is a second reason the physical space comes out three-dimensional, and it does not rely on the cusp at all.
Stable bound orbits, closed and neither precessing nor spiralling, exist only in three spatial dimensions. The
apsidal angle of the inverse-square law is exactly $180$ degrees in three dimensions. In four dimensions and
above, orbits either spiral inward or escape. This is the Ehrenfest and Bertrand result, and it constrains not
the substrate dimension but the dimension where matter lives. The natural reading is a braneworld, a
four-dimensional hyperbolic bulk of anti-de-Sitter character with matter confined to a three-dimensional slice.
That is the $\{3,4,3,4\}$ cusp picture exactly. So the claim that physical space is three-dimensional is doubly
forced, once by the cusp geometry and once by orbit stability, two arguments of different character landing on
the same number. The braneworld reading also carries a falsifiable handle, since a three-dimensional brane in a
four-dimensional bulk predicts short-range deviations from inverse-square gravity.


The dimension argument alone narrows the substrate to the four-dimensional family. It does not yet pick a single
honeycomb out of that family, and it does not need to, because the framework itself is dimension-general. The
rule, the Lorentz structure, electromagnetism, the holographic correlator, cosmology, and isotropy all port to
every dimension on the ladder. What the dimension does is set the stage. The crystallographic spinor coin,
treated next, is what selects $\{3,4,3,4\}$ within the four-dimensional family.
